% Preámbulo de los parciales · XeLaTeX · mismo estilo de los labs (rl_metodos/docs/curso/tex/preambulo.tex)
% Tipografía: Inter (cubre subíndices y griegas Unicode), con DejaVu Sans para los pocos símbolos que le faltan.
\documentclass[11pt,a4paper]{article}
\usepackage[margin=20mm,top=18mm,bottom=20mm]{geometry}
\usepackage{fontspec}
\setmainfont{Inter}[BoldFont={Inter SemiBold}, ItalicFont={Inter Italic}, BoldItalicFont={Inter SemiBold Italic}, Scale=0.95]
\setsansfont{Inter}
\setmonofont{DejaVu Sans Mono}[Scale=0.82]
\newfontfamily\dv{DejaVu Sans}[Scale=0.9]
\usepackage{unicode-math}
\setmathfont{latinmodern-math.otf}
\usepackage{polyglossia}\setdefaultlanguage{spanish}
\usepackage{xcolor,booktabs,longtable,array,calc,tabularx,enumitem,parskip,microtype,titlesec,fancyhdr,newunicodechar}
\usepackage[most]{tcolorbox}
\usepackage[hidelinks]{hyperref}
\emergencystretch=3em
\providecommand{\tightlist}{\setlength{\itemsep}{1pt}\setlength{\parskip}{0pt}}
\setcounter{secnumdepth}{-1}

\newunicodechar{∈}{{\dv ∈}}
\newunicodechar{∎}{{\dv ∎}}
\newunicodechar{⟨}{{\dv ⟨}}
\newunicodechar{⟩}{{\dv ⟩}}
\newunicodechar{☐}{{\dv ☐}}
\newunicodechar{✔}{{\dv ✔}}
\newunicodechar{✘}{{\dv ✘}}
\AtBeginDocument{\def\square{\mbox{\dv ☐}}}

\definecolor{tinta}{HTML}{1F2933}
\definecolor{acento}{HTML}{B7410E}
\definecolor{suave}{HTML}{F4F1EC}
\definecolor{gris}{HTML}{6B7280}
\definecolor{linea}{HTML}{D6D0C6}
\color{tinta}

\titleformat{\section}{\Large\bfseries\color{acento}}{}{0em}{}
\titleformat{\subsection}{\large\bfseries}{}{0em}{}
\titleformat{\subsubsection}{\normalsize\bfseries}{}{0em}{}
\titlespacing*{\section}{0pt}{1.4em}{0.4em}
\titlespacing*{\subsection}{0pt}{1.0em}{0.25em}
\setlist{nosep, leftmargin=1.4em, itemsep=2pt}
\renewcommand{\arraystretch}{1.3}
\setlength{\LTpre}{6pt}\setlength{\LTpost}{6pt}

% citas: el escenario inventado y las notas
\renewenvironment{quote}{\begin{tcolorbox}[colback=suave, colframe=linea, boxrule=0.5pt, arc=2pt, left=6pt, right=6pt, top=4pt, bottom=4pt]}{\end{tcolorbox}}
% código corto
\let\oldverbatim\verbatim
\let\endoldverbatim\endverbatim

% espacios para escribir
\newcommand{\hueco}[1]{\rule[-1.5pt]{\dimexpr #1\dimexpr 0.33em\relax\relax}{0.45pt}}
\newcommand{\lineasrespuesta}[1]{\par\vspace{2pt}{\color{linea}\foreach\i in {1,...,#1}{\vspace{0.75cm}\hrule height 0.5pt}}\par\vspace{4pt}}
\newcommand{\filarespuesta}{\rule{0pt}{1.25cm}}

\pagestyle{fancy}\fancyhf{}
\renewcommand{\headrulewidth}{0pt}
\fancyfoot[C]{\small\color{gris}\docpie\ · página \thepage}

\newcommand{\cabecera}[3]{%
  {\huge\bfseries\color{acento}#1\par}\vspace{4pt}
  {\large #2\par}\vspace{2pt}
  {\small\color{gris}#3\par}\vspace{6pt}
  {\color{linea}\hrule height 0.6pt}\vspace{8pt}}

\newcommand{\docpie}{Parcial 1 · Los dos Hermite}
\begin{document}
\cabecera{Los dos Hermite}{Polinomios osculantes, cuadratura de Gauss–Hermite y la molécula de HCl}{Entrega y sustentación: martes 20 de octubre de 2026 · Prof. Carlos Orrego}
En física, ``Hermite'' casi siempre quiere decir los polinomios Hₙ del oscilador armónico (H₀ = 1, H₁ = 2x, H₂ = 4x² − 2, \ldots). En este curso, ``Hermite'' es el polinomio osculante, el que copia valores y derivadas. Tienen el mismo apellido pero son objetos distintos, y en este parcial vas a ver que se necesitan el uno al otro: \textbf{la regla de cuadratura de Gauss--Hermite, que usan químicos y físicos para calcular valores esperados en el oscilador, se demuestra con el osculante}. Después la vas a usar en una molécula real, el cloruro de hidrógeno (HCl), cuyas constantes vibracionales medidas vienen de la tabla de NIST.

\section{Reglas del parcial}

Se entrega el \textbf{martes 20 de octubre al empezar la clase}, escrito a mano en esta misma guía, y se sustenta ese mismo día.

\begin{itemize}
\tightlist
\item
  \textbf{Individual.} Puedes discutir ideas con compañeros, pero lo que entregas lo escribes tú, con tus palabras.
\item
  \textbf{El cuaderno de Colab} de este parcial es \texttt{parcial\_hermite\_oscilador.ipynb}. Ábrelo desde el repositorio del curso, \textbf{github.com/cmorregof/methods\_exam}, carpeta \texttt{hermite-hcl}, con el botón «Open in Colab», y corre las celdas en orden. Antes de cada celda, escribe tu predicción en esta guía: la predicción escrita después de correr no cuenta.
\item
  \textbf{Puedes usar IA, y te audito.} Si la usas, llena la bitácora del final: qué le preguntaste, qué te respondió y qué corregiste o comprobaste tú. Usar IA no baja la nota. Lo que la baja es entregar algo que no sabes explicar.
\item
  \textbf{Sustentación (martes, 10 minutos).} Se sortea un punto de tu parcial y lo explicas en el tablero, sin la guía en la mano, con una variante (otro dato, otro valor inicial, otro nodo). La nota escrita se multiplica por un factor: 1 si lo explicas y lo reproduces; 0,6 si lo explicas a medias; 0,2 si no puedes explicar lo que entregaste.
\item
  \textbf{Puntaje.} Parte I: 12 puntos. Parte II: 14. Parte III: 12. Parte IV: 12. Total: 50. La nota es puntos ÷ 10.
\item
  \textbf{Qué cuenta.} Una respuesta correcta sin justificación vale poco. Una justificación clara con un error de cuenta vale mucho. Si algo te parece mal planteado, dilo y explica por qué: eso también se califica.
\end{itemize}

\section{Parte I · Completar una demostración (12 puntos)}

\textbf{Hechos que puedes usar sin demostrar.}

\begin{itemize}
\tightlist
\item
  (H1) Hₙ tiene n raíces reales y distintas x₁ \textless{} x₂ \textless{} \ldots{} \textless{} xₙ, y su coeficiente principal es 2ⁿ.
\item
  (H2) Ortogonalidad: si p es un polinomio de grado \textbf{menor que n}, entonces \ensuremath{\int} e\textsuperscript{−x²} Hₙ(x) p(x) dx = 0 (las integrales van de −∞ a ∞).
\item
  (H3) Osculante (forma de Burden, vista en clase). Sean Lⱼ los polinomios de Lagrange en los nodos x₁, \ldots, xₙ. El único polinomio de grado ≤ 2n − 1 con H(xⱼ) = f(xⱼ) y H′(xⱼ) = f′(xⱼ) es
\end{itemize}

\[H(x)=\sum_{j=1}^{n} f(x_j)\,\underbrace{\big[1-2(x-x_j)L_j'(x_j)\big]L_j(x)^2}_{H_j(x)}\;+\;\sum_{j=1}^{n} f'(x_j)\,\underbrace{(x-x_j)L_j(x)^2}_{\hat H_j(x)}\]

\[f(x)-H(x)=\frac{f^{(2n)}(\xi(x))}{(2n)!}\prod_{k=1}^{n}(x-x_k)^2 .\]

\begin{itemize}
\tightlist
\item
  (H4) Valor medio para integrales: si u(x) ≥ 0 y φ es continua, entonces \ensuremath{\int} φ(x) u(x) dx = φ(η) \ensuremath{\int} u(x) dx para algún η.
\end{itemize}

\textbf{Teorema (Gauss--Hermite).} Sean x₁, \ldots, xₙ las raíces de Hₙ. Para todo polinomio f de grado ≤ 2n − 1,

\[\int_{-\infty}^{\infty} e^{-x^2} f(x)\,dx=\sum_{j=1}^{n} w_j\,f(x_j)\]

\[w_j=\int_{-\infty}^{\infty} e^{-x^2}L_j(x)^2\,dx>0 .\]

Y si f tiene 2n derivadas continuas, el error es

\[\int_{-\infty}^{\infty} e^{-x^2} f(x)\,dx-\sum_{j=1}^{n} w_j f(x_j)=\frac{f^{(2n)}(\eta)}{(2n)!}\int_{-\infty}^{\infty} e^{-x^2}\prod_{k=1}^{n}(x-x_k)^2\,dx .\]

\subsection{I.1 · Llena los espacios (7 puntos)}

\textbf{Paso 1.} Multiplica el osculante por e\textsuperscript{−x²} e integra. Queda

\[\int e^{-x^2}H(x)\,dx=\sum_{j} f(x_j)\,A_j+\sum_{j} f'(x_j)\,B_j,\qquad A_j=\int e^{-x^2}H_j(x)\,dx,\quad B_j=\int e^{-x^2}(x-x_j)L_j(x)^2\,dx .\]

\textbf{Paso 2 (la clave).} Fija j. El polinomio q(x) = (x − xⱼ) Lⱼ(x) tiene grado \hueco{8} y se anula en \hueco{24}. Hₙ tiene ese mismo grado y esas mismas raíces, así que q(x) = c · Hₙ(x) para alguna constante c.~Para hallar c, deriva y evalúa en xⱼ:

q′(xⱼ) = Lⱼ(xⱼ) + (xⱼ − xⱼ) Lⱼ′(xⱼ) = \hueco{8}, y por tanto c = \hueco{16}.

\textbf{Paso 3.} Entonces

\[B_j=\int e^{-x^2}\,\big[(x-x_j)L_j(x)\big]\,L_j(x)\,dx = c\int e^{-x^2}H_n(x)\,\underline{\hspace{2cm}}\,dx=\underline{\hspace{1.5cm}}\]

por el hecho (H\hueco{4}), porque el grado de Lⱼ es \hueco{8}, que es menor que n.

\textbf{Paso 4.} Como Hⱼ = Lⱼ² − 2Lⱼ′(xⱼ)(x − xⱼ)Lⱼ², se tiene Aⱼ = \ensuremath{\int} e\textsuperscript{−x²} Lⱼ² dx − 2Lⱼ′(xⱼ) · \hueco{8} = wⱼ. Además wⱼ \textgreater{} 0 porque \hueco{40}.

\textbf{Paso 5 (exactitud).} Si f es un polinomio de grado ≤ 2n − 1, entonces H = f por \hueco{24} (cita H3). Por tanto \ensuremath{\int} e\textsuperscript{−x²} f dx = \ensuremath{\int} e\textsuperscript{−x²} H dx = Σ wⱼ f(xⱼ). \textbf{Las derivadas f′(xⱼ) desaparecieron: la regla solo necesita valores.}

\textbf{Paso 6 (error).} Restando, \ensuremath{\int} e\textsuperscript{−x²}(f − H) dx = \ensuremath{\int} {[}f\textsuperscript{(2n)}(ξ(x))/(2n)!{]} · e\textsuperscript{−x²} \ensuremath{\prod}(x − xₖ)² dx. Como la función \hueco{24} es ≥ 0, el hecho (H\hueco{4}) da el error del teorema.

\textbf{Remate.} Por (H1), \ensuremath{\prod}(x − xₖ) = Hₙ(x) / \hueco{8}. Usando el dato \ensuremath{\int} e\textsuperscript{−x²} Hₙ(x)² dx = 2ⁿ n! √π, el error queda

\[\text{error}=\frac{f^{(2n)}(\eta)}{(2n)!}\cdot\frac{n!\,\sqrt{\pi}}{\underline{\hspace{1.2cm}}} .\]

\subsection{I.2 · Pruébalo a mano con n = 2 (3 puntos)}

En el estado base del oscilador (ħ = m = ω = 1), la densidad de probabilidad es \textbar ψ₀(x)\textbar² = e\textsuperscript{−x²}/√π, así que ⟨f⟩ = (1/√π) \ensuremath{\int} e\textsuperscript{−x²} f(x) dx. Datos: \ensuremath{\int} e\textsuperscript{−x²} dx = √π, \ensuremath{\int} x² e\textsuperscript{−x²} dx = √π/2, y el valor exacto es ⟨x⁴⟩ = 3/4.

\begin{enumerate}
\def\labelenumi{\alph{enumi})}
\item
  Con H₂ = 4x² − 2, halla los nodos x₁, x₂ y calcula w₁ = \ensuremath{\int} e\textsuperscript{−x²} L₁(x)² dx. Escribe L₁ explícitamente. (Por simetría, w₂ = w₁.)
\item
  Calcula la aproximación de ⟨x⁴⟩ con la regla de 2 nodos.
\item
  Calcula el error con la fórmula del teorema (aquí f⁽⁴⁾ es constante, así que no necesitas η). ¿Aproximación + error = 3/4?
\end{enumerate}

\subsection{I.3 · ¿Dónde se usó\ldots? (2 puntos)}

\begin{enumerate}
\def\labelenumi{\alph{enumi})}
\item
  ¿En qué paso se usó que los nodos son \textbf{las raíces de Hₙ} y no unos nodos cualesquiera? Si tomas los nodos −1, 0, 1, ¿qué pasa con los Bⱼ? (Vas a ver la respuesta numérica en la Parte III, acción C.)
\item
  ¿En qué paso hace falta que e\textsuperscript{−x²} \ensuremath{\prod}(x − xₖ)² no cambie de signo? ¿Qué se perdería si cambiara?
\end{enumerate}

\begin{center}\rule{0.5\linewidth}{0.5pt}\end{center}

\section{Parte II · Predecir → correr → explicar con datos reales (14 puntos)}

\textbf{Los datos.} NIST Chemistry WebBook, página \emph{Hydrogen chloride}, tabla \emph{Constants of diatomic molecules} (compilación de Huber y Herzberg, 1979), estado fundamental X ¹Σ⁺ de H³⁵Cl: ωe = 2990,9463 cm⁻¹, ωe·xe = 52,8186 cm⁻¹, ωe·ye = 0,22437 cm⁻¹, ωe·ze = −0,01218 cm⁻¹, rₑ = 1,274552 Å. URL: https://webbook.nist.gov/cgi/cbook.cgi?ID=C7647010\&Mask=1000 (consultada el 15 oct 2026; el cuaderno trae una copia congelada por si la página falla).

\textbf{El modelo (te lo regalamos).} Mide la distancia en unidades de ℓ = √(ħ/(μω)) y la energía en unidades de ħω (= ωe). Entonces el potencial armónico es V\_arm(x) = x²/2 y el de Morse es

\[V_M(x)=D\,\big(1-e^{-a x}\big)^2,\qquad D=\frac{\omega_e}{4\,\omega_e x_e}\approx 14{,}157,\qquad a=\frac{1}{\sqrt{2D}}\approx 0{,}18793 .\]

(Con a = 1/√(2D), V\textsubscript{M} y V\_arm tienen la misma curvatura en x = 0.) La diferencia es g(x) = V\textsubscript{M}(x) − x²/2.

\textbf{Receta 1 (primer orden).} Un término extra W(x) corre el nivel base en ⟨ψ₀\textbar W\textbar ψ₀⟩ = (1/√π) \ensuremath{\int} e\textsuperscript{−x²} W(x) dx.

\textbf{Receta 2 (segundo orden, solo para un término cúbico b·x³).} Lo corre en −(11/8) b².

En cada ítem: \textbf{(a) predice en esta guía antes de correr}, \textbf{(b) corre la celda}, \textbf{(c) explica} la diferencia entre lo que predijiste y lo que salió.

\subsection{II.1 · La regla y ⟨x⁴⟩ --- Celdas 3, 4 y 5 (4 puntos)}

\begin{enumerate}
\def\labelenumi{\alph{enumi})}
\tightlist
\item
  \textbf{Predice.} En la Celda 4, ¿cuál de las tres columnas de pesos debería dar 0, según el Paso 3? ¿Cuál es el menor n con el que la regla da ⟨x⁴⟩ = 3/4 exacto, y qué da con n = 2?
\end{enumerate}

\lineasrespuesta{3}

\begin{enumerate}
\def\labelenumi{\alph{enumi})}
\setcounter{enumi}{1}
\item
  \textbf{Corre.} Completa la única línea de la Celda 3 (la regla (1/√π) Σ wⱼ g(xⱼ)) y corre las Celdas 3, 4 y 5.
\item
  \textbf{Explica.} ¿Por qué la columna que debería ser 0 da números como 10⁻¹⁷ y no 0? Cambia k = 6 en la Celda 5: ¿qué n mínimo aparece ahora y por qué? Cambia k = 5: ¿por qué sale exacto incluso con n = 1?
\end{enumerate}

\lineasrespuesta{3}

\lineasrespuesta{3}

\subsection{II.2 · La corrección de Morse con n nodos --- Celda 6 (4 puntos)}

La forma cerrada existe porque en la gaussiana ⟨e\textsuperscript{−kx}⟩ = e\textsuperscript{k²/4}: ⟨g⟩ = D(1 − 2e\textsuperscript{a²/4} + e\textsuperscript{a²}) − 1/4.

\begin{enumerate}
\def\labelenumi{\alph{enumi})}
\tightlist
\item
  \textbf{Predice.} e\textsuperscript{−ax} es una función entera (todas sus derivadas existen y son acotadas cerca de 0). Usa la fórmula del error de la Parte I con η ≈ 0 y el dato g\textsuperscript{(2n)}(0) = D a\textsuperscript{2n}(4ⁿ − 2) para n ≥ 2: al pasar de n a n + 1 nodos, ¿el error se divide por 2, por 10, por 100 o por más? ¿Con cuántos nodos crees que se llega al límite de la máquina (ε ≈ 2,2 × 10⁻¹⁶)?
\end{enumerate}

\lineasrespuesta{3}

\begin{enumerate}
\def\labelenumi{\alph{enumi})}
\setcounter{enumi}{1}
\item
  \textbf{Corre} la Celda 6. Mira las dos columnas de error y la gráfica.
\item
  \textbf{Explica.} ¿Qué tan cerca está la estimación de la Parte I del error real? Desde n = 8 las dos columnas de error se separan: con el valor ``ingenuo'' el error se estanca en unos 3 × 10⁻¹⁵, y con el ``estable'' baja hasta el orden de 10⁻¹⁶. Las dos fórmulas son algebraicamente idénticas. ¿Cuál de los temas del curso explica la diferencia? Señala la resta culpable.
\end{enumerate}

\lineasrespuesta{3}

\lineasrespuesta{3}

\subsection{II.3 · ¿Cuánto se corre el nivel base? --- Celda 7 (3 puntos)}

Taylor de Morse alrededor de 0 (usa D·a² = 1/2): V\textsubscript{M}(x) = x²/2 − (a/2) x³ + (7/24) a² x⁴ + \ldots{}

\begin{enumerate}
\def\labelenumi{\alph{enumi})}
\tightlist
\item
  \textbf{Predice.} Según la receta 1, ¿el nivel base sube o baja? Según la fórmula exacta de Morse, E\textsubscript{v} = ωe(v + ½) − ωe·xe(v + ½)², ¿cuántos cm⁻¹ se corre E₀ respecto al armónico?
\end{enumerate}

\lineasrespuesta{3}

\begin{enumerate}
\def\labelenumi{\alph{enumi})}
\setcounter{enumi}{1}
\item
  \textbf{Corre} la Celda 7, con \texttt{incluir\_cubico\ =\ False} y después con \texttt{True}.
\item
  \textbf{Explica.} ¿Por qué la receta 1 y Morse exacto dan signos contrarios? ¿Por qué el término cúbico no aporta nada a primer orden? (Piensa en paridad: ¿cuánto vale ⟨x³⟩?) ¿Qué pasa cuando sumas la receta 2?
\end{enumerate}

\lineasrespuesta{3}

\lineasrespuesta{3}

\subsection{II.4 · La frecuencia fundamental de HCl --- Celda 8 (3 puntos)}

La fundamental es la transición v = 0 → 1: ν = E₁ − E₀. Para el armónico, ν = ωe; para Morse, ν = ωe − 2ωe·xe. La ``observada'' se calcula con todas las constantes de NIST, que vienen de ajustar las bandas infrarrojas medidas.

\begin{enumerate}
\def\labelenumi{\alph{enumi})}
\tightlist
\item
  \textbf{Predice.} ¿Qué error relativo (en \%) esperas del modelo armónico? ¿Y del de Morse? ¿Cuántos niveles ligados tiene el pozo de Morse? (Pista: los niveles existen mientras E\textsubscript{v} siga creciendo con v.)
\end{enumerate}

\lineasrespuesta{3}

\begin{enumerate}
\def\labelenumi{\alph{enumi})}
\setcounter{enumi}{1}
\item
  \textbf{Corre} la Celda 8.
\item
  \textbf{Explica.} ¿Por qué el armónico \textbf{sobreestima} la frecuencia? Mira en la tabla cómo se separan E\_arm y E\_Morse al subir v. ¿A qué longitud de onda (en μm) absorbe HCl su fundamental?
\end{enumerate}

\lineasrespuesta{3}

\lineasrespuesta{3}

\begin{center}\rule{0.5\linewidth}{0.5pt}\end{center}

\section{Parte III · El agente (12 puntos)}

\begin{quote}
\textbf{Escenario inventado para este parcial.} Los números de la tabla sí están calculados de verdad, pero el agente no existe.
\end{quote}

Un agente de RL tiene un \textbf{presupuesto de 6 condiciones} para estimar valores esperados en el estado base, ⟨f⟩ = (1/√π) \ensuremath{\int} e\textsuperscript{−x²} f(x) dx. Una condición es un valor f(xⱼ) o una derivada f′(xⱼ). En cada acción escoge dónde poner las condiciones y usa la regla que integra exactamente el polinomio (Lagrange u osculante) que pasa por ellas. Su recompensa es −\textbar error\textbar{} en ⟨g⟩ = ⟨V\textsubscript{M} − V\_arm⟩ de HCl (valor exacto 0,0078276501608889). Las acciones E y F están como referencia.

\begin{longtable}[]{@{}
  >{\raggedright\arraybackslash}p{(\columnwidth - 10\tabcolsep) * \real{0.1667}}
  >{\raggedright\arraybackslash}p{(\columnwidth - 10\tabcolsep) * \real{0.1667}}
  >{\raggedright\arraybackslash}p{(\columnwidth - 10\tabcolsep) * \real{0.1667}}
  >{\raggedright\arraybackslash}p{(\columnwidth - 10\tabcolsep) * \real{0.1667}}
  >{\raggedright\arraybackslash}p{(\columnwidth - 10\tabcolsep) * \real{0.1667}}
  >{\raggedright\arraybackslash}p{(\columnwidth - 10\tabcolsep) * \real{0.1667}}@{}}
\toprule\noalign{}
\begin{minipage}[b]{\linewidth}\raggedright
Acción
\end{minipage} & \begin{minipage}[b]{\linewidth}\raggedright
Condiciones
\end{minipage} & \begin{minipage}[b]{\linewidth}\raggedright
Pesos de las derivadas
\end{minipage} & \begin{minipage}[b]{\linewidth}\raggedright
⟨x⁴⟩ (exacto 0,75)
\end{minipage} & \begin{minipage}[b]{\linewidth}\raggedright
⟨x⁶⟩ (exacto 1,875)
\end{minipage} & \begin{minipage}[b]{\linewidth}\raggedright
\textbar error\textbar{} en ⟨g⟩ de HCl
\end{minipage} \\
\midrule\noalign{}
\endhead
\bottomrule\noalign{}
\endlastfoot
A · Gauss--Hermite n = 6 & 6 valores & --- & 0,75 & 1,875 & 2,7 × 10⁻¹² \\
B · osculante en los 3 nodos de H₃ & 3 valores + 3 derivadas & 0, 0, 0 & 0,75 & 1,125 & 4,1 × 10⁻⁵ \\
C · osculante en −1, 0, 1 & 3 valores + 3 derivadas & −0,0625; 0; 0,0625 & 0,75 & 1,000 & 4,8 × 10⁻⁵ \\
D · 6 nodos equiespaciados en {[}−2,5; 2,5{]} & 6 valores & --- & 0,75 & 1,984 & 5,9 × 10⁻⁶ \\
E · Gauss--Hermite n = 3 & 3 valores & --- & 0,75 & 1,125 & 4,1 × 10⁻⁵ \\
F · Gauss--Hermite n = 2 & 2 valores & --- & 0,25 & 0,125 & 5,2 × 10⁻³ \\
\end{longtable}

\textbf{El agente aprendió a escoger A.}

\subsection{III.1 · ¿Verdadero o falso? (6 puntos, 1 cada una)}

Para cada afirmación del ``reporte del agente'', escribe V o F y \textbf{justifícala con la teoría} en una o dos líneas. Un V o F sin razón no vale.

\begin{enumerate}
\def\labelenumi{\arabic{enumi}.}
\tightlist
\item
  ``Con n = 2 nodos de Gauss--Hermite, la regla da exactamente ⟨x⁵⟩, aunque 5 es mayor que 2n − 1 = 3.''
\item
  ``Con n = 2 nodos la regla es exacta para todo polinomio de grado ≤ 5.''
\item
  ``La acción B gasta 6 condiciones y rinde exactamente lo mismo que E con 3: sus 3 derivadas no aportan nada.''
\item
  ``B y C usan 6 condiciones cada una, así que son exactas para todo polinomio de grado ≤ 2·6 − 1 = 11.''
\item
  ``Pasar de 3 a 6 nodos de Gauss--Hermite (de E a A) baja el error de ⟨g⟩ más o menos a la mitad.''
\item
  ``En C las derivadas sí cuentan (sus pesos no son todos 0), porque −1, 0, 1 no son las raíces de H₃.''
\end{enumerate}

\subsection{III.2 · Inferencia (3 puntos)}

¿Por qué el agente aprendió A y no B, si las dos gastan 6 condiciones? Responde con el número 2n − 1 y con el Paso 3 de la Parte I. ¿Y por qué A le gana a D, que también usa 6 valores?

\lineasrespuesta{3}

\lineasrespuesta{3}

\subsection{III.3 · Reward hacking (3 puntos)}

Una versión anterior del agente se evaluaba sobre la familia \{x³, x⁵, sin x\}, en vez de ⟨g⟩, y pagaba un pequeño costo por cada condición usada. Ese agente aprendió a usar \textbf{un solo nodo} (Gauss--Hermite n = 1, el nodo x = 0) y reportó ``error 0 en todas las pruebas''.

\begin{enumerate}
\def\labelenumi{\alph{enumi})}
\tightlist
\item
  ¿Por qué obtuvo error 0 de verdad? b) ¿Qué error tiene esa acción en ⟨x⁴⟩ y en ⟨g⟩ de HCl? (Calcúlalo: no hace falta el computador.) c) Propón un cambio concreto a la familia de prueba para que esta trampa deje de funcionar.
\end{enumerate}

\lineasrespuesta{3}

\lineasrespuesta{3}

\begin{center}\rule{0.5\linewidth}{0.5pt}\end{center}

\section{Parte IV · Consulta: nuestra idea de paper aplicada a tu tema (12 puntos)}

Haz una búsqueda profunda (NotebookLM Deep Research u otra herramienta) sobre \textbf{cuadratura de Gauss y Gauss--Hermite en física y química cuántica}: la representación de variable discreta (DVR), los métodos espectrales y pseudoespectrales para la ecuación de Schrödinger, y trabajos de ML/RL que escogen bases, mallas o nodos en problemas cuánticos.

Algunos puntos de partida que el profesor verificó que existen. Igual tienes que auditarlos, porque que existan no quiere decir que digan lo que alguien les atribuye:

\begin{itemize}
\tightlist
\item
  G. H. Golub y J. H. Welsch, ``Calculation of Gauss quadrature rules'', \emph{Math. Comp.} 23 (1969), 221--230.
\item
  J. C. Light, I. P. Hamilton y J. V. Lill, ``Generalized discrete variable approximation in quantum mechanics'', \emph{J. Chem. Phys.} 82 (1985), 1400--1409.
\item
  D. T. Colbert y W. H. Miller, ``A novel discrete variable representation for quantum mechanical reactive scattering via the S-matrix Kohn method'', \emph{J. Chem. Phys.} 96 (1992), 1982--1991.
\item
  J. Yang et al., ``Reinforcement Learning for Adaptive Mesh Refinement'', arXiv:2103.01342 (2021).
\end{itemize}

\subsection{IV.1 · Síntesis (4 puntos)}

Una página escrita \textbf{por ti, sin IA}, que responda las preguntas guía:

\begin{enumerate}
\def\labelenumi{\arabic{enumi}.}
\tightlist
\item
  ¿Qué es una DVR y qué tiene que ver con los nodos de Gauss y los Lⱼ de la Parte I? (Pista: ¿qué le pasa a la matriz del potencial cuando la calculas con la regla de cuadratura?)
\item
  ¿Por qué los métodos espectrales convergen ``exponencialmente'' para funciones suaves, y cómo se relaciona con lo que viste en la Celda 6?
\item
  ¿Hay trabajos donde un agente de RL o un modelo de ML escoja la base, la malla o los nodos en un problema cuántico? ¿Qué usan como recompensa? ¿Tienen riesgo de reward hacking como el de III.3?
\end{enumerate}

\subsection{IV.2 · Cuatro fuentes auditadas (4 puntos, 1 por fuente)}

Al menos \textbf{dos} tienen que ser fuentes que encontraste tú, no de la lista de arriba.

\begin{longtable}[]{@{}
  >{\raggedright\arraybackslash}p{(\columnwidth - 10\tabcolsep) * \real{0.1667}}
  >{\raggedright\arraybackslash}p{(\columnwidth - 10\tabcolsep) * \real{0.1667}}
  >{\raggedright\arraybackslash}p{(\columnwidth - 10\tabcolsep) * \real{0.1667}}
  >{\raggedright\arraybackslash}p{(\columnwidth - 10\tabcolsep) * \real{0.1667}}
  >{\raggedright\arraybackslash}p{(\columnwidth - 10\tabcolsep) * \real{0.1667}}
  >{\raggedright\arraybackslash}p{(\columnwidth - 10\tabcolsep) * \real{0.1667}}@{}}
\toprule\noalign{}
\begin{minipage}[b]{\linewidth}\raggedright
\#
\end{minipage} & \begin{minipage}[b]{\linewidth}\raggedright
Referencia (autores, año, título)
\end{minipage} & \begin{minipage}[b]{\linewidth}\raggedright
¿Existe? (DOI/enlace)
\end{minipage} & \begin{minipage}[b]{\linewidth}\raggedright
Tipo (artículo, libro, preprint\ldots)
\end{minipage} & \begin{minipage}[b]{\linewidth}\raggedright
¿Dice lo que la herramienta le atribuye?
\end{minipage} & \begin{minipage}[b]{\linewidth}\raggedright
Conexión con el curso
\end{minipage} \\
\midrule\noalign{}
\endhead
\bottomrule\noalign{}
\endlastfoot
1 \filarespuesta & & & & & \\
2 \filarespuesta & & & & & \\
3 \filarespuesta & & & & & \\
4 \filarespuesta & & & & & \\
\end{longtable}

\subsection{IV.3 · Nuestra idea en tu tema (4 puntos)}

La idea de paper del curso es una biblioteca de demostraciones escritas como grafo (pasos, razones, dependencias, distractores), sobre la que un agente aprende a demostrar, más agentes que toman las mismas decisiones que un estudiante de métodos numéricos.

\begin{enumerate}
\def\labelenumi{\alph{enumi})}
\item
  (2 puntos) Escribe la demostración de la Parte I como grafo: una lista de 6 a 8 nodos (paso + razón), con flechas ``depende de'', e incluye \textbf{un distractor}: un paso que suene plausible pero sea falso (por ejemplo, una razón equivocada para que Bⱼ = 0).
\item
  (2 puntos) \textbf{¿Qué tendría que ser verdad para que nuestra idea fuera nueva en tu tema?} Piensa en un agente que aprenda a escoger nodos o bases para problemas cuánticos y que además ``sepa por qué'' (que pueda citar el Paso 3 o la fórmula del error). ¿Qué encontraste, o no encontraste, en tu búsqueda que lo haga nuevo o ya hecho?
\end{enumerate}

\lineasrespuesta{3}

\lineasrespuesta{3}

\section{Bitácora de IA}

Una fila por cada uso. Si no usaste IA, escribe «No usé IA» y firma.

\begin{longtable}[]{@{}
  >{\raggedright\arraybackslash}p{(\columnwidth - 6\tabcolsep) * \real{0.1047}}
  >{\raggedright\arraybackslash}p{(\columnwidth - 6\tabcolsep) * \real{0.2674}}
  >{\raggedright\arraybackslash}p{(\columnwidth - 6\tabcolsep) * \real{0.3140}}
  >{\raggedright\arraybackslash}p{(\columnwidth - 6\tabcolsep) * \real{0.3140}}@{}}
\toprule\noalign{}
\begin{minipage}[b]{\linewidth}\raggedright
Punto
\end{minipage} & \begin{minipage}[b]{\linewidth}\raggedright
Qué le pedí
\end{minipage} & \begin{minipage}[b]{\linewidth}\raggedright
Qué me respondió, en una línea
\end{minipage} & \begin{minipage}[b]{\linewidth}\raggedright
Qué comprobé o corregí yo, y cómo
\end{minipage} \\
\midrule\noalign{}
\endhead
\bottomrule\noalign{}
\endlastfoot
\filarespuesta & & & \\
\filarespuesta & & & \\
\filarespuesta & & & \\
\filarespuesta & & & \\
\filarespuesta & & & \\
\filarespuesta & & & \\
\filarespuesta & & & \\
\end{longtable}

\textbf{Firma:} \hueco{30}

\section{Sustentación (la llena el profesor el martes)}

\textbf{Punto sorteado:} \hueco{14}

\begin{itemize}
\tightlist
\item[$\square$]
  Explica el punto con sus palabras, sin leer la guía
\item[$\square$]
  Reproduce en el tablero la cuenta o el paso clave
\item[$\square$]
  Responde una variante: otro dato, otro valor inicial, otro nodo
\item[$\square$]
  Lo que dice coincide con lo que entregó y con su bitácora
\end{itemize}

\textbf{Factor:} 1 ☐ 0,6 ☐ 0,2 ☐ \textbf{Nota escrita:} \hueco{6} / 50 \textbf{Nota final:} \hueco{6}

\end{document}
